% GENERATED FILE. Edit canonical lessons, then run npm run generate:books.
% canonical-source-hash: fnv1a64:91ac7265fceb6f3d
% canonical-lessons: HI-C470-sahas, HI-C470-vir, HI-C470-diler, HI-C470-nirbhik, HI-C470-jokhim

\chapter{Courage}
\label{ch:hi-courage}
\hlchaptermodality{470}

\begin{chapteropening}
\textbf{By the end of this chapter:} \emph{I can say that someone was brave.}

Complete the last lesson of chapter 470: I can say that someone was brave.
\end{chapteropening}

\section[\texorpdfstring{\hi{साहस} (sāhas)}{साहस (sāhas)}]{\hi{साहस} (sāhas) — daring; boldness}

\label{lesson:HI-C470-sahas}



\begin{quote}
Before the new one: say the Hindi for greedy, then the Hindi for a beggar.
\end{quote}

\subsection*{You'll want to know: \hi{साहस}}
\textbf{\hi{साहस}} — \emph{sāhas} — \textquotedblleft{}daring; boldness\textquotedblright{}.

\begin{quote}
\hi{रात} \hi{में} \hi{अकेले} \hi{जंगल} \hi{में} \hi{जाने} \hi{के} \hi{लिए} \hi{साहस} \hi{चाहिए।}
\end{quote}

\begin{quote}
\emph{rāt meṁ akele jaṅgal meṁ jāne ke liye sāhas cāhiye.}
\end{quote}

\begin{quote}
It takes daring to go into the forest alone at night.
\end{quote}

Masculine.

\begin{cousinweb}
\textbf{\hi{साहस}} (\emph{sāhas}) is a tatsama, borrowed from Sanskrit \emph{sāhasa}, \textquotedblleft{}boldness, daring\textquotedblright{}, also \textquotedblleft{}rashness\textquotedblright{} (Monier-Williams).
\end{cousinweb}

\begin{grammarlens}[title={What you've built}]
One more word for describing what you lived and felt.
\end{grammarlens}

\subsection*{Your turn}
\begin{itemize}
\raggedright
  \item \emph{Say it:} \emph{sāhas}
  \item \emph{Say it:} \emph{sāhas}, once more
  \item \emph{Say it:} \emph{rāt meṁ akele jaṅgal meṁ jāne ke liye sāhas cāhiye.} — It takes daring to go into the forest alone at night.
  \item \emph{Recall:} say the Hindi for greedy, then the Hindi for a beggar, then say \emph{sāhas} again
\end{itemize}

\begin{tcolorbox}[breakable,colback=teal!4,colframe=teal!35!black,title={Before you move on}]
What does \hi{साहस} mean? (\textquotedblleft{}Daring; boldness\textquotedblright{}.) Say it once more.
\end{tcolorbox}
\section[\texorpdfstring{\hi{वीर} (vīr)}{वीर (vīr)}]{\hi{वीर} (vīr) — valiant; a hero}

\label{lesson:HI-C470-vir}



\begin{quote}
Before the new one: say the Hindi for a beggar, then the Hindi for daring.
\end{quote}

\subsection*{You'll want to know: \hi{वीर}}
\textbf{\hi{वीर}} — \emph{vīr} — \textquotedblleft{}valiant; a hero\textquotedblright{}.

\begin{quote}
\hi{दादा} \hi{एक} \hi{वीर} \hi{सिपाही} \hi{थे।}
\end{quote}

\begin{quote}
\emph{dādā ek vīr sipāhī the.}
\end{quote}

\begin{quote}
Grandfather was a valiant soldier.
\end{quote}

\begin{cousinweb}
\textbf{\hi{वीर}} (\emph{vīr}) is a tatsama, borrowed from Sanskrit \emph{vīra}, \textquotedblleft{}a man, esp. a brave one; a hero\textquotedblright{} (Monier-Williams).
\end{cousinweb}

\begin{grammarlens}[title={What you've built}]
One more word for describing what you lived and felt.
\end{grammarlens}

\subsection*{Your turn}
\begin{itemize}
\raggedright
  \item \emph{Say it:} \emph{vīr}
  \item \emph{Say it:} \emph{vīr}, once more
  \item \emph{Say it:} \emph{dādā ek vīr sipāhī the.} — Grandfather was a valiant soldier.
  \item \emph{Recall:} say the Hindi for a beggar, then the Hindi for daring, then say \emph{vīr} again
\end{itemize}

\begin{tcolorbox}[breakable,colback=teal!4,colframe=teal!35!black,title={Before you move on}]
What does \hi{वीर} mean? (\textquotedblleft{}Valiant; a hero\textquotedblright{}.) Say it once more.
\end{tcolorbox}
\section[\texorpdfstring{\hi{दिलेर} (diler)}{दिलेर (diler)}]{\hi{दिलेर} (diler) — bold; fearless}

\label{lesson:HI-C470-diler}



\begin{quote}
Before the new one: say the Hindi for daring, then the Hindi for valiant.
\end{quote}

\subsection*{You'll want to know: \hi{दिलेर}}
\textbf{\hi{दिलेर}} — \emph{diler} — \textquotedblleft{}bold; fearless\textquotedblright{}.

\begin{quote}
\hi{मीरा} \hi{बहुत} \hi{दिलेर} \hi{है}, \hi{चोर} \hi{को} \hi{देखकर} \hi{भी} \hi{नहीं} \hi{डरी।}
\end{quote}

\begin{quote}
\emph{Mīrā bahut diler hai, cor ko dekhkar bhī nahīṃ ḍarī.}
\end{quote}

\begin{quote}
Mira is very bold; she wasn't afraid even when she saw the thief.
\end{quote}

\begin{cousinweb}
\textbf{\hi{दिलेर}} (\emph{diler}) is usually traced to Persian \emph{dilēr}, \textquotedblleft{}bold\textquotedblright{}, built on \emph{dil}, \textquotedblleft{}heart\textquotedblright{}.
\end{cousinweb}

\begin{grammarlens}[title={What you've built}]
One more word for describing what you lived and felt.
\end{grammarlens}

\subsection*{Your turn}
\begin{itemize}
\raggedright
  \item \emph{Say it:} \emph{diler}
  \item \emph{Say it:} \emph{diler}, once more
  \item \emph{Say it:} \emph{Mīrā bahut diler hai, cor ko dekhkar bhī nahīṃ ḍarī.} — Mira is very bold; she wasn't afraid even when she saw the thief.
  \item \emph{Recall:} say the Hindi for daring, then the Hindi for valiant, then say \emph{diler} again
\end{itemize}

\begin{tcolorbox}[breakable,colback=teal!4,colframe=teal!35!black,title={Before you move on}]
What does \hi{दिलेर} mean? (\textquotedblleft{}Bold; fearless\textquotedblright{}.) Say it once more.
\end{tcolorbox}
\section[\texorpdfstring{\hi{निर्भीक} (nirbhīk)}{निर्भीक (nirbhīk)}]{\hi{निर्भीक} (nirbhīk) — fearless}

\label{lesson:HI-C470-nirbhik}



\begin{quote}
Before the new one: say the Hindi for valiant, then the Hindi for bold.
\end{quote}

\subsection*{You'll want to know: \hi{निर्भीक}}
\textbf{\hi{निर्भीक}} — \emph{nirbhīk} — \textquotedblleft{}fearless\textquotedblright{}.

\begin{quote}
\hi{निर्भीक} \hi{पत्रकार} \hi{ने} \hi{सच} \hi{लिखा।}
\end{quote}

\begin{quote}
\emph{nirbhīk patrakār ne sac likhā.}
\end{quote}

\begin{quote}
The fearless journalist wrote the truth.
\end{quote}

\begin{cousinweb}
\textbf{\hi{निर्भीक}} (\emph{nirbhīk}) is a tatsama, borrowed from Sanskrit \emph{nirbhīka}, \textquotedblleft{}fearless\textquotedblright{}, built on \emph{niḥ}, \textquotedblleft{}without\textquotedblright{}, and \emph{bhī}, \textquotedblleft{}fear\textquotedblright{} (Monier-Williams), the root that also gives \emph{bhay}.
\end{cousinweb}

\begin{grammarlens}[title={What you've built}]
One more word for describing what you lived and felt.
\end{grammarlens}

\subsection*{Your turn}
\begin{itemize}
\raggedright
  \item \emph{Say it:} \emph{nirbhīk}
  \item \emph{Say it:} \emph{nirbhīk}, once more
  \item \emph{Say it:} \emph{nirbhīk patrakār ne sac likhā.} — The fearless journalist wrote the truth.
  \item \emph{Recall:} say the Hindi for valiant, then the Hindi for bold, then say \emph{nirbhīk} again
\end{itemize}

\begin{tcolorbox}[breakable,colback=teal!4,colframe=teal!35!black,title={Before you move on}]
What does \hi{निर्भीक} mean? (\textquotedblleft{}Fearless\textquotedblright{}.) Say it once more.
\end{tcolorbox}
\section[\texorpdfstring{\hi{जोखिम} (jokhim)}{जोखिम (jokhim)}]{\hi{जोखिम} (jokhim) — a risk}

\label{lesson:HI-C470-jokhim}



\begin{quote}
Before the new one: say the Hindi for bold, then the Hindi for fearless.
\end{quote}

\subsection*{You'll want to know: \hi{जोखिम}}
\textbf{\hi{जोखिम}} — \emph{jokhim} — \textquotedblleft{}a risk\textquotedblright{}.

\begin{quote}
\hi{पहाड़} \hi{पर} \hi{अकेले} \hi{चढ़ने} \hi{में} \hi{बहुत} \hi{जोखिम} \hi{है।}
\end{quote}

\begin{quote}
\emph{pahāṛ par akele caṛhne meṁ bahut jokhim hai.}
\end{quote}

\begin{quote}
Climbing the mountain alone is very risky.
\end{quote}

Masculine.

\begin{cousinweb}
\textbf{\hi{जोखिम}} (\emph{jokhim}) is a tadbhava: CDIAL lists Hindi forms of this word, with the senses \textquotedblleft{}charge for securing property from accidents, risk, danger, loss\textquotedblright{}, under Sanskrit \emph{yōgakṣēmá}, \textquotedblleft{}preservation of one's property; an insurance charge\textquotedblright{} (CDIAL 10527): the price of safety became the danger itself.
\end{cousinweb}

\begin{grammarlens}[title={What you've built}]
One more word for describing what you lived and felt.
\end{grammarlens}

\subsection*{Your turn}
\begin{itemize}
\raggedright
  \item \emph{Say it:} \emph{jokhim}
  \item \emph{Say it:} \emph{jokhim}, once more
  \item \emph{Say it:} \emph{pahāṛ par akele caṛhne meṁ bahut jokhim hai.} — Climbing the mountain alone is very risky.
  \item \emph{Recall:} say the Hindi for bold, then the Hindi for fearless, then say \emph{jokhim} again
\end{itemize}

\begin{tcolorbox}[breakable,colback=teal!4,colframe=teal!35!black,title={Before you move on}]
What does \hi{जोखिम} mean? (\textquotedblleft{}A risk\textquotedblright{}.) Say it once more.
\end{tcolorbox}
